% ==================================================
% EXPERIMENTS & RESULTS
% ==================================================

\section{Experiments}

\subsection{Experimental Setup}

\subsubsection{Dataset Statistics}

The current EDA was generated by the pipeline in \texttt{src/mars\_rec/eda}. It loads the MARS interaction and metadata files, validates the required schema, merges implicit and explicit interactions, filters active interactions, computes dataset statistics, and produces plots in \texttt{outputs/eda/figures}.

\begin{table}[H]
    \centering
    \caption{MARS active interaction statistics from EDA}
    \label{tab:eda-statistics}
    \begin{tabular}{lr}
        \toprule
        \textbf{Statistic} & \textbf{Value} \\
        \midrule
        Users & 3,007 \\
        Items & 958 \\
        Active interactions & 24,622 \\
        Sparsity & 99.1453\% \\
        Density & 0.8547\% \\
        Mean interactions per user & 8.19 \\
        Min/max interactions per user & 1 / 395 \\
        Mean users per item & 25.70 \\
        Min/max users per item & 1 / 886 \\
        Users with sequence length $\geq 3$ & 1,676 \\
        \bottomrule
    \end{tabular}
\end{table}

\begin{table}[H]
    \centering
    \caption{Item metadata coverage}
    \label{tab:metadata-coverage}
    \begin{tabular}{lr}
        \toprule
        \textbf{Metadata field} & \textbf{Coverage} \\
        \midrule
        Difficulty & 40.70\% \\
        Software & 100.00\% \\
        Theme & 100.00\% \\
        Duration & 100.00\% \\
        Type & 100.00\% \\
        Description & 86.46\% \\
        \bottomrule
    \end{tabular}
\end{table}

\begin{figure}[H]
    \centering
    \includegraphics[width=0.82\textwidth]{../outputs/eda/figures/sequence_length_distribution.png}
    \caption{Sequence length distribution in the active interaction dataset}
    \label{fig:sequence-length-distribution}
\end{figure}

\begin{figure}[H]
    \centering
    \includegraphics[width=0.82\textwidth]{../outputs/eda/figures/item_popularity_long_tail.png}
    \caption{Long-tail item popularity pattern}
    \label{fig:item-popularity-long-tail}
\end{figure}

\subsubsection{Implementation Details}

The planned implementation stack is:

\begin{itemize}
    \item \textbf{Language:} Python.
    \item \textbf{Model framework:} PyTorch.
    \item \textbf{Data processing:} pandas and parquet intermediate files.
    \item \textbf{Demo:} Streamlit with a Python/FastAPI backend at prototype scale.
    \item \textbf{Database:} SQLite for lightweight demo storage.
\end{itemize}

\subsubsection{Hyperparameter Tuning}

We conducted grid search over:

\begin{itemize}
    \item Embedding dimensions: \{32, 64, 128\}
    \item Number of attention heads: \{1, 2, 4, 8\}
    \item Number of layers: \{1, 2, 3\}
    \item Learning rates: \{0.001, 0.0005, 0.0001\}
\end{itemize}

Optimal hyperparameters were selected based on validation set performance.
[TBD: Replace this sentence with the final selected hyperparameters and validation criterion after the training runs are completed.]

\subsection{Results}

\subsubsection{Main Performance Results}

\begin{table}[htbp]
    \centering
    \caption{Sequential Recommendation Performance Comparison}
    \label{tab:performance}
    \begin{tabular}{|c|c|c|c|c|}
        \hline
        \textbf{Method} & \textbf{Recall@5} & \textbf{NDCG@5} & \textbf{MRR@5} & \textbf{Hit Rate@5} \\
        \hline
        Popularity & [TBD] & [TBD] & [TBD] & [TBD] \\
        \hline
        Item-based CF & [TBD] & [TBD] & [TBD] & [TBD] \\
        \hline
        BPR-MF & [TBD] & [TBD] & [TBD] & [TBD] \\
        \hline
        SASRec & [TBD] & [TBD] & [TBD] & [TBD] \\
        \hline
        \textbf{gSASRec Direction} & \textbf{[TBD]} & \textbf{[TBD]} & \textbf{[TBD]} & \textbf{[TBD]} \\
        \hline
    \end{tabular}
\end{table}

\noindent [TBD: Insert final quantitative results after model training, validation, and test evaluation.]

\subsubsection{Ablation Study}

To understand the contribution of different components, the planned ablation study will compare the following variants:

\begin{table}[htbp]
    \centering
    \caption{Ablation Study Results}
    \label{tab:ablation}
    \begin{tabular}{|c|c|}
        \hline
        \textbf{Model Variant} & \textbf{Recall@5} \\
        \hline
        gSASRec direction (full setting) & [TBD] \\
        \hline
        Standard BCE instead of Generalised BCE & [TBD] \\
        \hline
        Single negative sample instead of multiple negatives & [TBD] \\
        \hline
        Without metadata initialization & [TBD] \\
        \hline
        Without confidence weighting & [TBD] \\
        \hline
        Without diversity/novelty reranking & [TBD] \\
        \hline
    \end{tabular}
\end{table}

\subsection{Analysis}

\subsubsection{Performance Insights}

[TBD: To be filled with detailed analysis of model performance, including which model performs best under Recall@K, nDCG@K, and HitRate@K.]

\subsubsection{Learning Dynamics}

[TBD: To be filled with training curves, convergence behavior, validation performance, and early-stopping observations.]

\subsubsection{Error Analysis}

The expected error analysis will examine:

\begin{itemize}
    \item Cold-start or low-history users.
    \item Rare resources in the long tail.
    \item Cases where popularity dominates personalized signals.
    \item Sequences with unusual or non-linear learning patterns.
\end{itemize}

\subsection{Statistical Significance}

[TBD: Add significance testing or confidence intervals if repeated runs are performed. Otherwise, state the residual evaluation risk clearly.]

\newpage
